\chapter*{Основные обозначения и сокращения}
\addcontentsline{toc}{chapter}{Основные обозначения и сокращения}

\section*{Сокращения}

\begin{description}
	\item[ЭРДУ] электрореактивная двигательная установка.
	\item[КА] космический аппарат.
	\item[ИР] идеально регулируемый.
	\item[ПСИ] постоянная скорость истечения.
	\item[РИД] результат интеллектуальной деятельности.
	\item[КЛО] кривая локальных оптимумов.
\end{description}

\section*{Основные обозначения и термины}

\begin{description}
	\item[$\mathbf{r}$] радиус-вектор положения КА.
	\item[$r$] длина радиус-вектора положения КА.
	\item[$\mathbf{v}$] вектор скорости КА.
	\item[$\mathbf{a}$] реактивное ускорение.
	\item[$\mathbf{p}$] вектор сопряжённых переменных.
	\item[$T$] продолжительность перелёта.
	\item[$J$] значение квадратичного функционала.
	\item[$\kappa$] отношение радиусов орбит.
	\item[$\varphi$] угловая дальность.
	\item[$\lambda$] витковая дальность.
	\item[$\Phi$] синодическая фаза.
	\item[$R$] матрица линейного преобразования, или $R$-преобразование.
	\item[$i$] наклонение конечной орбиты.
	\item[$\Omega$] долгота восходящего узла.
	\item[$\bar{\omega}$] долгота перицентра.
	\item[$\omega$] аргумент перицентра.
	\item[DU] единица длины в принятой безразмерной системе.
	\item[TU] единица времени в принятой безразмерной системе.
	\item[модельное многообразие] специально построенное многообразие решений энергетически оптимальных задач встречи в упрощённой постановке.
	\item[аппроксимированное многообразие] аналитический объект, полученный после сборки конечных формул для характеристик траектории и параметров управления.
	\item[точки оптимальной дальности] точки пересечения КЛО с изолиниями начальной синодической фазы.
	\item[линейные преобразования] общий класс преобразований, используемых для перехода от плоской задачи к пространственной и эллиптической постановке.
	\item[$R$-преобразование] специальное линейное преобразование параметров решения, используемое для перехода от плоской задачи к пространственной и эллиптической постановке.
	\item[начальное приближение] набор параметров решения, получаемый по аппроксимациям и линейным преобразованиям до последующей оптимизации.
\end{description}